\begin{helper}
For any finite set of ordered pairs of base vertices, normalizing each ordered
pair to the increasing coordinate can identify at most two ordered pairs with
one normalized coordinate.  Consequently, the number of ordered pairs is at
most twice the number of normalized coordinates.
\end{helper}
\begin{proof}
Let \(N(x,y)\) be \((x,y)\) if \(x<y\), and \((y,x)\) otherwise.  Fix a
normalized coordinate \(q\).  If \(N(e)=q\), then either \(e=q\), in the
increasing case, or \(e\) is the reversal of \(q\), in the other case.  Thus
every fiber of the normalization map has cardinality at most two.  Applying
the finite fiber-cardinality bound \(\noderef{ProjectionFiberCardBound}\) gives
the desired inequality.
\end{proof}
